\chapter{Thực nghiệm và kết quả}
\label{ch:thucnghiem}

\section{Thiết lập thực nghiệm}

Ba nhánh được chấm trên toàn bộ quần thể bước chạm của tập kiểm: mô hình gốc chưa tinh
chỉnh (\textbf{Base}), mô hình sau tinh chỉnh trên nhánh chỉ sinh câu, hạt giống $101$
(\textbf{S1}), và câu do người viết (\textbf{Trần}). Trong $4.463$ bước, S1 trả về câu
rỗng ở đúng một bước, tính là trượt. Mọi điểm tuyệt đối và khoảng tin cậy trong chương
này đo trên đủ $4.463$ bước; riêng các phép so ghép cặp chạy trên $4.462$ bước được cả ba
nhánh trả lời. Hai quần thể chênh nhau một bước nên không đại lượng nào đổi sau khi làm
tròn tới một chữ số thập phân.

Khâu sinh câu chạy trên card A100 ngay sau khi huấn luyện, khi trọng số còn nằm sẵn trong
bộ nhớ. Khâu chấm chạy trên card T4 miễn phí, mất khoảng $5{,}6$ giờ cho mỗi nhánh, và
không tốn chi phí.

\subsection{Ghi nhận về lượt huấn luyện}

Lượt S1 hạt giống $101$ chạy đủ $8.072$ bước tương ứng hai lượt duyệt toàn bộ $64.567$
mẫu, tệp bộ chuyển đổi đầu ra nặng $59{,}9$ MB. Tổng số phép tính \code{total\_flos} là
$4.142.257.957$ GF, tức $513.164$ GF cho mỗi bước.

Con số đó là bằng chứng cho thấy toàn bộ $8.072$ bước đã thực sự chạy, dù lượt huấn luyện
bị gián đoạn sáu lần do mất máy ảo, và cơ chế chạy tiếp từ điểm lưu không làm mất cũng
không làm lặp bước nào. Cách lập luận ở đây đáng nói, vì bản trước của chương này dùng một
lập luận \emph{vòng tròn}: nó so \code{total\_flos} với chính phép ngoại suy từ lượt thăm
dò đã đẻ ra kỳ vọng ấy. Bằng chứng thay thế không có khuyết điểm đó là so hai lượt huấn
luyện \emph{độc lập} có số lần gián đoạn khác hẳn nhau --- hạt giống $101$ đứt sáu lần,
hạt giống $202$ đứt ba lần --- mà \code{total\_flos} trên mỗi bước khớp nhau trong
$0{,}013\%$ ($513.164$ so với $513.229$). Nếu khâu kế toán sai sau mỗi lần chạy tiếp thì
hai lượt phải lệch nhau theo số lần đứt.

Cũng cần ghi nhận rằng ngoại suy từ lượt thăm dò lệch tới $7{,}7\%$ chứ không phải
$0{,}32\%$, vì lượt thăm dò chạy với \code{max\_samples} nên chỉ lấy phần đầu của tập
\emph{chưa trộn}, độ dài chuỗi không đại diện cho toàn tập.

Đường mất mát gộp theo $200$ bước đi từ $1{,}103$ ở giai đoạn đầu, giai đoạn mô hình học
định dạng đầu ra, xuống $0{,}611$ ở bước $2.000$, có một nhịp tụt đúng tại ranh giới lượt
duyệt thứ hai ở bước $4.036$, và phẳng ở mức $0{,}455$ về cuối. Đuôi phẳng là do tốc độ
học đã xuống $2{,}28 \times 10^{-7}$ theo đúng lịch cosine; đó là dấu hiệu kết thúc sạch
chứ không phải dấu hiệu hỏng. Cần lưu ý hai điều: nhịp tụt ở lượt duyệt thứ hai \emph{không} phải bằng
chứng khái quát tốt hơn, vì gặp lại dữ liệu lần hai thì mất mát huấn luyện giảm là đương
nhiên; và không nên ngoại suy đường mất mát: một dự báo sàn từng được đưa ra rồi phải
rút vì thực tế đi thấp hơn nhiều.

Một ghi chú về tái lập: sau khi chạy tiếp từ điểm lưu, ba trường trong tệp kết quả tổng
hợp của khung huấn luyện bị sai và không được trích dẫn. Trường mất mát huấn luyện báo
$0{,}0151$ vì khung cộng mất mát chỉ từ lúc khởi động lại ($272$ bước) rồi chia cho toàn
bộ $8.072$ bước; giá trị thật xấp xỉ $0{,}447$, kiểm được bằng phép nhân ngược. Mọi đại
lượng mà khung chia cho thời gian trôi qua đều hỏng theo cùng cơ chế. Các số đúng lấy từ
tệp trạng thái của bộ huấn luyện.

\section{Kết quả chính}

\begin{table}[t]
\caption{Thước đo áp lên ba nhánh trên quần thể $4.463$ bước chạm, kèm hai thước dựa trên
câu chuẩn. Executability đọc trên nền $75{,}7\%$.}
\label{tab:chinh}
\centering\small
\renewcommand{\arraystretch}{1.25}
\begin{tabular}{@{}lccrr@{}}
\toprule
Nhánh & Exec. & KTC 95\% & BLEU-4 & ROUGE-L \\
\midrule
Base (chưa tinh chỉnh) & $47{,}6$ & $[45{,}9; 49{,}3]$ & $9{,}7$ & $40{,}3$ \\
S1 (hạt giống $101$) & $59{,}1$ & $[57{,}3; 60{,}8]$ & $38{,}7$ & $67{,}3$ \\
S1 (hạt giống $202$) & $59{,}6$ & $[57{,}9; 61{,}3]$ & --- & --- \\
S1 (trung bình hai hạt giống) & $\mathbf{59{,}4}$ & --- & --- & --- \\
\midrule
\emph{Câu do người viết (trần)} & $75{,}7$ & $[74{,}1; 77{,}3]$ & $96{,}1$ & $100{,}0$ \\
\bottomrule
\end{tabular}

\vspace{0.4em}
\begin{minipage}{0.95\columnwidth}\footnotesize
Ba nhánh còn lại (S2 hai hạt giống, S2r, S2-nopoint) cùng nhánh
B-infer đã đăng ký và đã dựng xong dữ liệu nhưng chưa chạy xong tại thời điểm nộp;
trạng thái từng nhánh ở Bảng~\ref{tab:trangthai}.
\end{minipage}
\end{table}

Bảng~\ref{tab:chinh} cho kết quả chính. Bảng~\ref{tab:phanra} tách điểm thành các thành
phần hợp thành để thấy phần nào đóng góp vào đâu.

\begin{table}[t]
\caption{Phân rã điểm số theo từng điều kiện của thước.}
\label{tab:phanra}
\centering\small
\renewcommand{\arraystretch}{1.2}
\begin{tabular}{@{}lrrrr@{}}
\toprule
Nhánh & Executable & Chỉ ô Voronoi & Khớp thao tác & Ngưỡng dung sai \\
\midrule
Trần (câu người viết) & $75{,}7$ & $75{,}8$ & $100{,}0$\textsuperscript{$*$} & $84{,}3$ \\
S1 (hạt giống $101$) & $59{,}1$ & $60{,}2$ & $94{,}4$ & $69{,}2$ \\
S1 (hạt giống $202$) & $59{,}6$ & $60{,}7$ & $94{,}6$ & $69{,}8$ \\
Base & $47{,}6$ & $48{,}9$ & $96{,}5$ & $57{,}6$ \\
\bottomrule
\multicolumn{5}{@{}p{12.4cm}@{}}{\footnotesize \textsuperscript{$*$}Vì câu dự đoán chính
là câu chuẩn nên điều kiện khớp thao tác đúng theo định nghĩa; điểm rút gọn về phần định
vị. Đó chính là ý nghĩa của trần.}
\end{tabular}
\end{table}

\paragraph{Một kết cục có thể chấm dứt nghiên cứu đã bị loại trừ.} Nếu nhánh nền đã sát
mức câu người viết thì không còn chỗ cho bất kỳ can thiệp nào, và bản đăng ký trước yêu
cầu dừng lại trong trường hợp đó. Điều đó không xảy ra: ba khoảng tin cậy rời nhau hoàn
toàn. Thước cũng không suy biến về sàn.

\section{So sánh ghép cặp}

Vì cả ba nhánh được chấm trên cùng một tập bước, phép so đúng là phép so ghép cặp
(McNemar). Bảng~\ref{tab:mcnemar} cho kết quả; cả ba đều có $p < 0{,}001$.

\begin{table}[t]
\caption{So sánh ghép cặp trên $4.462$ bước. $b$ là số bước nhánh dưới giải được mà nhánh
trên không giải được, $c$ là chiều ngược lại.}
\label{tab:mcnemar}
\centering\small
\renewcommand{\arraystretch}{1.2}
\begin{tabular}{@{}lrrrr@{}}
\toprule
Phép so & $\Delta$ & $b$ & $c$ & $\chi^2$ \\
\midrule
S1 $-$ Base & $+11{,}5$ & $284$ & $798$ & $243{,}2$ \\
Trần $-$ S1 & $+16{,}6$ & $102$ & $843$ & $579{,}5$ \\
Trần $-$ Base & $+28{,}1$ & $102$ & $1.357$ & $1.077{,}8$ \\
\bottomrule
\end{tabular}
\end{table}

Dòng giữa là dòng đáng chú ý nhất về mặt thiết kế: có \textbf{$843$ bước mà câu do người
viết định vị được còn S1 thì không}. Đó chính là vùng mà nhánh S2 có thể cải thiện, và nó
là dạng ghép cặp của con số ``dư địa'' ở Mục~\ref{sec:dudia}.

Sai số chuẩn của hiệu ghép cặp là $0{,}79$ điểm dưới bootstrap gom cụm, so với $1{,}05$
điểm nếu tính như hai tỉ lệ độc lập.

Với lượt huấn luyện này, tinh chỉnh trên tập giám sát dựng tự động đáng giá $11{,}5$ điểm
và đóng được $41\%$ khoảng cách giữa mô hình chưa tinh chỉnh và mức câu người viết. Tỉ lệ
đó \emph{cao hơn thực chất} do một phần tư số bước mà S1 chép lại câu chuẩn: trên $3.367$ bước
mà mô hình tự viết ra câu của mình, ba nhánh đọc lần lượt $44{,}7$, $53{,}1$ và $75{,}0$,
và phần đóng được chỉ là $28\%$.

\paragraph{Mức câu người viết không phải chặn trên của cái đạt được.} Hợp của ba nhánh
giải được $3.527$ bước, tức $79{,}1\%$, và trong số các bước mà câu người viết trượt có
$148$ bước được một mô hình giải đúng. Cái mà $75{,}7\%$ đo là một lớp câu cụ thể đi được
tới đâu với dụng cụ này.

\section{Lặp lại lượt huấn luyện: nhiễu của cả đường ống}
\label{sec:haihatgiong}

Một chênh lệch chỉ đáng tin nếu nó lớn hơn mức mà \emph{dựng lại chính hệ thống đó} làm
điểm xê dịch. Luận văn huấn luyện S1 lần thứ hai với hạt giống khác, mọi khoá cấu hình còn
lại kiểm khớp $38/38$.

\begin{table}[t]
\caption{Hai lượt huấn luyện chỉ khác hạt giống, chấm trên cùng $4.462$ bước.}
\label{tab:haihatgiong}
\centering\small
\renewcommand{\arraystretch}{1.2}
\begin{tabular}{@{}lrrr@{}}
\toprule
& Điểm & $\Delta$ ghép cặp & Bước bất đồng \\
\midrule
S1 hạt giống $101$ & $59{,}12\%$ & \multirow{2}{*}{$+0{,}52$ pp} & \multirow{2}{*}{$287 = 6{,}4\%$} \\
S1 hạt giống $202$ & $59{,}64\%$ & & \\
\midrule
\emph{đối chiếu}: S1 $-$ Base & & $+11{,}5$ pp & $1.082 = 24{,}2\%$ \\
\bottomrule
\end{tabular}
\end{table}

KTC 95\% của chênh lệch giữa hai hạt giống là $[-0{,}21; +1{,}28]$ --- \emph{chứa số không},
và McNemar không bác ($b = 132$, $c = 155$, $p = 0{,}194$). Đó đúng là thứ một giả thuyết
không phải cho ra.

\paragraph{Tỉ lệ bước bất đồng phân biệt tốt hơn con số chênh lệch.} Hai lượt chỉ khác hạt
giống bất đồng ở $6{,}4\%$ số bước; một can thiệp thật làm bất đồng $24{,}2\%$. Chênh lệch
$11{,}5$ điểm vì vậy lớn gấp \textbf{$22$ lần} nhiễu dựng lại về độ lớn, và gấp $3{,}8$ lần
về tỉ lệ xáo trộn.

\paragraph{Không một phần nào của nhiễu ấy đến từ thước.} Trên $2.810$ bước mà hai lượt
viết ra câu \emph{trùng nhau từng byte}, mô hình định vị trả về toạ độ giống hệt và phán
quyết giống hệt, $2.810/2.810$. Toàn bộ biến thiên quan sát được là biến thiên của
\emph{hệ thống được đo}, không phải của \emph{dụng cụ đo}.

\paragraph{Độ tin cậy lặp lại, và một con số dễ đọc nhầm.} Mức đồng thuận từng bước giữa
hai lượt là $93{,}6\%$ ($\kappa = 0{,}867$). Nhưng $63\%$ số cặp là câu trùng từng byte,
nơi một thước tất định \emph{buộc} phải đồng thuận. Trên $1.652$ bước hai lượt thật sự viết
khác nhau, con số là $82{,}6\%$ ($\kappa = 0{,}650$), và \textbf{đó mới là số nên trích}.
Điểm theo từng ứng dụng tương quan $r = 0{,}927$ trên $18$ ứng dụng có ít nhất $20$ bước.

\paragraph{Mọi lát cắt chẩn đoán đều lặp lại.} Sáu lát ở Mục~\ref{sec:latcat} được tính
lại trên hạt giống thứ hai: cả sáu cùng dấu, lệch nhau nhiều nhất $0{,}9$ điểm. Riêng lát
``ứng dụng chưa thấy lúc dạy'' trùng khít, $+14{,}1$ và $+14{,}1$. Thắng--hoà--thua theo cụm
cũng vậy: $395/565/131$ và $406/572/113$ trên $1.091$ cụm.

\section{Kết luận này có đứng vững không?}
\label{sec:phanbien}

Chênh lệch $11{,}5$ điểm cho phép nhiều cách đọc khác nhau. Bảy cách đọc thay thế được
xem xét trên bản ghi từng bước; năm cách bị bác, hai cách còn lại chưa bác được và được
giữ nguyên trong bảng thay vì bỏ đi (Bảng~\ref{tab:phanbien}).

\begin{table}[t]
\caption{Bảy cách giải thích thay thế cho chênh lệch S1 trên Base, và kết quả kiểm.
Năm cách bị bác; hai cách cuối chưa bác được.}
\label{tab:phanbien}
\centering\small
\renewcommand{\arraystretch}{1.25}
\begin{tabular}{@{}p{3.2cm}p{9.2cm}@{}}
\toprule
Cách đọc thay thế & Kết quả kiểm \\
\midrule
Do điều kiện khớp thao tác gánh &
bỏ hẳn điều kiện đó, chấm riêng phần định vị, vẫn còn $+11{,}3$ điểm ($b=272$, $c=777$).
Hơn nữa Base gọi \emph{đúng} loại thao tác nhiều hơn ($96{,}5\%$ so với $94{,}4\%$) \\[0.3em]

Chỉ thắng ở một vị trí bước &
bước mở đầu tác vụ $+15{,}0$ ($n=593$, $[+11{,}0; +19{,}1]$), các bước sau $+11{,}0$
($[+9{,}4; +12{,}6]$); hai khoảng giao nhau nên không kết luận được thứ tự theo vị trí,
nhưng cả hai đều dương \\[0.3em]

Chỉ thắng ở màn dễ &
chia ba theo số phần tử trên màn: $+12{,}6$, $+9{,}5$, $+12{,}4$; các khoảng giao nhau,
không tầng nào đảo chiều \\[0.3em]

Vài ứng dụng kéo cả bảng &
trên $1.091$ cụm, S1 thắng $395$, hoà $565$, thua $131$ \\[0.3em]

Chỉ là bắt chước văn phong &
\textbf{không bác được theo cách này}: S1 mở đầu bằng động từ chạm nhiều hơn Base
($91{,}0\%$ so với $86{,}1\%$; câu người viết $95{,}1\%$), tức nó thực sự nằm gần văn
phong người chú thích hơn \\[0.3em]

Do dụng cụ đo đã quen văn phong &
\textbf{chưa kiểm được}: đòi hỏi một mô hình định vị thứ hai đã xác minh sạch
AndroidControl. Đây là cách đọc thay thế mạnh nhất còn lại (Mục~\ref{sec:hanche}) \\[0.3em]

Do chép lại câu chuẩn &
$1.095$ trong $4.462$ câu của S1 lặp đúng các từ nội dung của câu chuẩn. Chấm riêng
$3.367$ bước mà mô hình tự viết, S1 vẫn hơn $+8{,}4$ điểm ($b=264$, $c=547$,
$\chi^2 = 98{,}1$, $p<0{,}001$) \\
\bottomrule
\end{tabular}
\end{table}

Dòng cuối là dòng quan trọng nhất và là dòng mà một thước dựa trên câu chuẩn không thể tự
kiểm giúp. Lợi thế của S1 lớn hơn hẳn trên nhóm bước chép lại ($+21{,}1$ điểm), nên một
phần những gì tinh chỉnh mua được đúng là cách diễn đạt nguyên văn của người chú thích;
nhưng thứ tự giữa hai nhánh không dựa vào phần đó.

Riêng lập luận về \emph{độ dài câu} cần bàn tách ra. Thước có thiên vị câu dài: chia đôi ở độ
dài trung vị $33$ ký tự cho $61{,}9\%$ so với $56{,}5\%$, và trên nhóm bước mà S1 tự viết
thì chênh lệch là $9{,}6$ điểm. Tuy nhiên, thiên vị đó \emph{nghiêng về phía Base}: độ dài
trung vị câu của Base là khoảng $70$ ký tự, của S1 là $33$, của người viết là $34$. S1
thắng \emph{dù chịu bất lợi} về độ dài, nên kết luận vì vậy mạnh lên. Cần nói
thêm rằng thiên vị này không phải quy luật chung của thước, vì trong nội bộ nhánh Base
câu dài lại kém hơn $1{,}4$ điểm, nên nó phải được đo riêng cho từng nhánh.

\section{Lỗi nằm ở đâu}
\label{sec:latcat}

Bốn quan sát từ bản ghi từng bước ràng buộc cách đọc nhánh xử lý về sau.

\subsection{Không quan sát thấy lợi thế sân nhà}

Đây là phản biện nặng nhất mà luận văn tự nêu ra: nếu $95{,}6\%$ số bước quy được tên của tập
kiểm thuộc về ứng dụng cũng có ở tập dạy thì điểm số bị thổi lên. Bảng~\ref{tab:sannha} đo trực tiếp.

\begin{table}[t]
\caption{Điểm S1 theo tình trạng ứng dụng của bước.}
\label{tab:sannha}
\centering\small
\begin{tabular}{@{}lrr@{}}
\toprule
Nhóm & S1 & $n$ \\
\midrule
Ứng dụng đã thấy lúc dạy & $59{,}1\%$ & $1.737$ \\
Ứng dụng \emph{chưa} thấy & $59{,}0\%$ & $78$ \\
Không quy được về ứng dụng & $59{,}2\%$ & $2.647$ \\
\bottomrule
\end{tabular}
\end{table}

Ba con số gần như trùng nhau. Với $78$ bước thì khoảng tin cậy của nhóm chưa thấy rất
rộng, nên phát biểu đúng là \emph{không phát hiện được ảnh hưởng}, chứ không phải \emph{đã
chặn được ảnh hưởng}. Dù vậy, đây là một giới hạn được nêu kèm số đo bên cạnh:
con số $95{,}6\%$ đọc rất nặng cho tới khi có bảng này đặt bên cạnh.

\subsection{Lỗi còn lại là lỗi mô tả, không phải lỗi chọn thao tác}

Mô hình gọi đúng loại thao tác ở $94{,}4\%$ số bước. Trong $1.824$ bước trượt, chỉ $251$
bước sai loại thao tác; $1.573$ bước còn lại, tức $86\%$, có loại thao tác đúng và một
cách diễn đạt mà mô hình định vị không lần ra được. Vài ví dụ thật:

\begin{center}\footnotesize\ttfamily
\setlength{\tabcolsep}{4pt}
\begin{tabular}{@{}ll@{}}
\toprule
\rmfamily câu chuẩn & \rmfamily câu S1 sinh ra \\
\midrule
click on the OK option & select minutes 00 \\
Click on the Sort filter & Click on the Air India flight option \\
Click on the share icon & Open Recorder app \\
\bottomrule
\end{tabular}
\end{center}

Phần lỗi còn lại nằm đúng ở chỗ \emph{cách gọi tên và mô tả phần tử}, tức đúng chỗ mà một
đích huấn luyện có tầng khai báo nhắm tới. Khác với phép thử tạo động lực ban đầu, quan
sát này được đo trên mô hình đã huấn luyện thật và trên toàn bộ quần thể.

\subsection{Tinh chỉnh thực sự dạy được gì}

Lấy từ $798$ bước mà S1 đúng còn Base sai, hai kiểu hỏng của Base hiện ra rõ:

\begin{center}\small
\renewcommand{\arraystretch}{1.15}
\begin{tabularx}{0.95\textwidth}{@{}lX@{}}
\toprule
chuẩn & \ttfamily Click on the profile icon at the top right of the screen \\
S1 & \ttfamily Click on the profile icon at the top right corner of the screen \\
Base & \ttfamily To set YouTube video downloads over Wi-Fi only, tap on your profile
icon at the top right corner of the screen and select ``Settings.'' Then navigate to the
``Downloads'' section and choose ``Wi-Fi Only.'' \\
\midrule
chuẩn & \ttfamily Click on Knoxville \\
S1 & \ttfamily Click on Knoxville \\
Base & \ttfamily Click on the three dots next to Knoxville to access more options and
proceed with reading the local news. \\
\bottomrule
\end{tabularx}
\end{center}

Kiểu thứ nhất là trả lời cả kế hoạch thay vì một bước; kiểu thứ hai là lạc sang một bước
khác. Cả hai đều không phải lỗi về khả năng viết tiếng Anh: Base viết câu trôi chảy,
đúng khuôn, gọi tên nút cụ thể. Thứ tinh chỉnh dạy được là \textbf{bám đúng bước hiện
tại}. Điều này làm cho Base trở thành một mốc so sạch, không lẫn với chuyện ``mô hình
không theo khuôn định dạng''.

\subsection{Dư địa nhỏ hơn số bước trượt gợi ý}
\label{sec:dudia}

Câu người viết cũng trượt ở $1.083$ trong số các bước này, và như Mục~\ref{sec:vungmu}
cho thấy, phần lớn là giới hạn của dụng cụ. Dư địa quy được về cách diễn đạt chưa tốt vào
khoảng $741$ bước, tức $16{,}6$ điểm; dạng ghép cặp của nó là $843$ bước mà câu người viết
giải được còn S1 thì không.

Đặt cạnh MDE ghép cặp đo được là $2{,}2$ điểm, một can thiệp chỉ cần sửa khoảng $13\%$ dư
địa đó là đủ để phát hiện. Mức hiệu ứng trung vị mà văn liệu báo cáo cho các can thiệp
tương đương nằm trên dải này, nhưng một kết quả không có tác dụng vẫn là kết cục còn sống
và đã có luật đọc sẵn.

\section{Vùng mù của thước}
\label{sec:vungmu}

Câu người viết trượt ở $1.083$ bước, và không phải tất cả đều do mô hình định vị.

Ở $64{,}8\%$ số bước đó, điểm trả về nằm hẳn ngoài vùng dung sai, tức mô hình định vị
\emph{bỏ cuộc} chứ không phải định vị lệch. Sai số của nó lưỡng cực rõ rệt: trung vị
$0{,}4\%$ ở những bước thành công so với $26{,}2\%$ ở những bước thất bại, gần như không
có vùng giữa.

$381$ bước còn lại vượt qua ngưỡng dung sai nhưng bị chính luật ô Voronoi bác; ở $208$
bước trong số đó, điểm nằm ngay bên trong hộp nhỏ nhất chứa điểm chạm; trên một phần tử
lớn, một phần tử con cạnh tranh với chính phần tử cha của nó. Tính theo nhánh, luật ô
Voronoi lật $10{,}1\%$ số bước đã qua ngưỡng dung sai của câu người viết, $13{,}0\%$ của S1
và $15{,}2\%$ của Base.

Theo chiều ngược lại, thước mù với các độ lệch dưới khoảng $63$ điểm ảnh
(Mục~\ref{sec:bomloi}). Và $935$ bước, tức $21\%$ quần thể, đánh bại cả ba nhánh.

Từ đó rút ra phát biểu quan trọng nhất về cách đọc mọi con số trong chương này:
\textbf{trần $75{,}7\%$ là giới hạn của dụng cụ đo, không phải giới hạn của ngôn ngữ}.
Nghi ngờ rằng trần thấp là do câu chuẩn quá cụt cũng đã được kiểm và bác: lọc riêng những
bước có câu chuẩn từ ba từ trở xuống cho trần $77{,}0\%$, không khá hơn bao nhiêu.

Một phép kiểm liên quan tới lo ngại ``mô hình có thể leo thước bằng cách viết cho dụng cụ
đo thay vì viết cho người'': câu nêu tên một vùng trên màn đạt $89{,}5\%$ khi vùng đó
đúng và $3{,}3\%$ khi vùng đó sai, so với $54{,}8\%$ cho câu không nêu vùng nào. Nghĩa là
phần thưởng dành cho việc \emph{nói đúng} vị trí, không phải cho việc \emph{có nhắc tới}
vị trí. S1 dùng từ chỉ vị trí ở $26{,}8\%$ số câu so với $26{,}4\%$ của người viết, và
chuỗi toạ độ chỉ xuất hiện ở $0{,}04\%$ số câu, đúng bằng tỉ lệ của người viết. Điều S1
làm nhiều hơn là dùng khuôn mẫu: hai mươi câu phổ biến nhất của nó phủ $10{,}4\%$ quần
thể so với $6{,}5\%$ của người viết. Nhánh S2, vốn được huấn luyện để sinh toạ độ, cần
được kiểm lại điểm này trước khi đọc điểm số của nó.

\section{Các nhánh đã đăng ký nhưng chưa có kết quả}
\label{sec:chuachay}

Tại thời điểm hoàn thành bản này, các nhánh sau đã đăng ký và đã dựng xong dữ liệu nhưng
chưa chạy xong. Các ô số tương ứng được để trống.

\begin{table}[t]
\caption{Trạng thái các nhánh còn lại và điều kiện phụ thuộc.}
\label{tab:trangthai}
\centering\small
\renewcommand{\arraystretch}{1.2}
\begin{tabular}{@{}p{2.9cm}p{3.4cm}p{5.9cm}@{}}
\toprule
Nhánh & Trạng thái & Phụ thuộc \\
\midrule
S1 hạt giống $202$ & \textbf{đã xong, đã chấm} & nhiễu hạt giống $+0{,}52$ điểm, Mục~\ref{sec:haihatgiong} \\
S2 $\times 2$ hạt giống & chưa chạy & sau khi khoá ngưỡng từ cặp hạt giống S1 \\
S2r & chưa chạy & cần để quy công cho tính phân biệt \\
S2-nopoint & chưa chạy & cần để tách phần công của ô toạ độ \\
B-infer & chưa chạy & không chặn nhánh nào \\
S3 (khoản phạt mức hai) & \textbf{chưa có mã hàm mất mát} & sau cổng riêng \\
Chấm tay $100$ câu, hai người & chưa chạy & cần câu do mô hình sinh \\
Kiểm không-gây-hại & chưa chạy & chạy cùng lượt chấm S2 \\
Kiểm chéo bằng mô hình định vị thứ hai & chưa chạy & xử lý giới hạn ở Mục~\ref{sec:hanche} \\
\bottomrule
\end{tabular}
\end{table}

Đại lượng chính của luận văn, $\Delta = \text{S2} - \text{S1}$, vì vậy \textbf{chưa được
báo cáo}, và luận văn không viết bất kỳ phát biểu nào về hiệu quả của thành phần ``mô tả
phân biệt trước, phát ngôn sau''. Việc thay thế đại lượng đó bằng phép so S1 với Base,
vốn đã có số, sẽ là đổi câu hỏi sau khi đã nhìn thấy dữ liệu, và bản đăng ký trước cấm
điều đó.

Khi các lượt chạy hoàn tất, các con số cần điền vào đúng ba chỗ: dòng tương ứng trong
Bảng~\ref{tab:chinh}, một bảng ghép cặp theo mẫu Bảng~\ref{tab:mcnemar}, và ba lát cắt đã
công bố trước ở Mục~\ref{sec:dangky}.
